% =====================================================================
%  27.4  Tarjan's offline LCA -- one DFS + DSU, WHITE/GREY/BLACK, union a
%  child on return, answer a query when its second endpoint turns BLACK.
%  O((n+q) alpha(n)) with NO lifting table; OFFLINE only.
%  WORKED EXAMPLE ONLY -- no solvedproblem (chapter contract: exactly two,
%  already spent in 27.1/27.2).  Back-cites Ch.16 (\chref{c11}) as the
%  ONLINE owner: lifting is used, never re-derived; here Tarjan is the
%  batch alternative.
%  Exports: sec:c21-tarjan  def:c21-tarjan  thm:c21-tarjan  lst:c21-tarjan
%  Trace runs on the authoritative R of 00_opener with the batch
%  Q1=(8,6)->1, Q2=(4,5)->2, Q3=(2,8)->2, Q4=(6,7)->3, Q5=(5,7)->1;
%  every union/answer event was hand-replayed against the DFS order
%  P=[1,2,4,8,5,3,6,7] before writing.
% =====================================================================
\section{Tarjan's Offline LCA: One DFS and a Disjoint-Set Union}\label{sec:c21-tarjan}

\begin{previouslybox}[Ch.~16 answers one LCA at a time -- this answers the batch]
Chapter~16 (\chref{c11}) built \(\lca\) from binary lifting: \(O(n\log n)\)
preprocessing, then each query climbs in \(O(\log n)\), \emph{online}.  We have
been using that oracle as a black box (\(\secref{c21-path}\) splits windows by
\(\lca\)).  Now suppose the \(q\) LCA queries are all known in advance -- can the
batch be cheaper than \(q\) separate climbs?  Tarjan's 1979 answer says yes, and
it is the most beautiful five-line idea in this chapter: run \emph{one plain
DFS}, keep a disjoint-set union of the region you have \emph{finished}, and
whenever a node dies (turns BLACK) it settles every old query whose other
endpoint already died.  Total \(\bigO{(n+q)\alpha(n)}\) -- effectively linear,
no lifting table, no \(\log\).
\end{previouslybox}

\begin{topicintro}
Paint nodes \emph{WHITE} (unseen), \emph{GREY} (on the DFS stack -- entered,
children pending), \emph{BLACK} (finished).  When a child \(v\) returns to its
parent \(u\), \(\mathit{union}\) \(v\)'s set into \(u\)'s and stamp
\(\mathit{ancestor}[\mathit{find}(u)] = u\): each DSU set is exactly one
maximal BLACK region, and its stamp is the GREY parent that region is waiting
under.  Attach every query \((u_i,v_i)\) to \emph{both} endpoints' lists.  When
\(u\) turns BLACK, scan its list: a partner already BLACK means \emph{this} is
the moment the two endpoints' paths met -- the answer is
\(\mathit{ancestor}[\mathit{find}(\text{partner})]\).  \thmref{c21-tarjan} proves
it and bounds it.  The worked trace runs five queries on the running tree \(R\),
including both sibling cases and a strict ancestor case.

\textbf{Prerequisites.} DSU with path compression is used as a black box here
(\(\mathit{find}/\mathit{union}\), amortised \(\alpha\)); \chref{c11} owns what
an \(\lca\) \emph{is} and the online lifting alternative this section is
measured against; the WHITE/GREY/BLACK DFS vocabulary is the standard
cycle-detection language from the traversal chapters.

\textbf{After this section you can}
\begin{itemize}[nosep]
  \item run the three-colour DFS with union-on-return and state the invariant; \hfill\lv{4}
  \item answer a batch of \(q\) LCAs in \(\bigO{(n+q)\alpha}\) offline; \hfill\lv{4}
  \item trace \(\mathit{ancestor}[\mathit{find}(\cdot)]\) answers by hand (done: five on \(R\)); \hfill\lv{5}
  \item choose between Tarjan, lifting, and Euler+RMQ from the statement's modality. \hfill\lv{4}
\end{itemize}
\end{topicintro}

% ---------------------------------------------------------------
\subsection{The algorithm and its invariant}
% ---------------------------------------------------------------

\begin{definition}{Tarjan's offline LCA}{c21-tarjan}
Root the tree; attach each query \((u_i,v_i)\) to both \(u_i\) and \(v_i\).
DFS: on \emph{entering} \(u\) set \(\mathit{col}[u]=\)GREY, \(\mathit{fa}[u]=u\),
\(\mathit{anc}[u]=u\); after each child \(v\) returns,
\(\mathit{fa}[\mathit{find}(v)] := \mathit{find}(u)\) and
\(\mathit{anc}[\mathit{find}(u)] := u\); when all children are done, for each
attached query whose partner is BLACK record
\(\mathit{ans}_i = \mathit{anc}[\mathit{find}(\text{partner})]\), then set
\(\mathit{col}[u]:=\)BLACK.
\end{definition}

\begin{theorem}{Tarjan answers every offline LCA in \(\bigO{(n+q)\alpha(n)}\)}{c21-tarjan}
At every instant, each DSU set is precisely a maximal fully-BLACK connected
region, and \(\mathit{anc}[\mathit{find}(x)]\) is the deepest GREY node whose
subtree contains \(x\).  When a query's second endpoint \(u\) turns BLACK,
that stamp equals \(\lca(u,v)\).  Total time \(\bigO{(n+q)\alpha(n)}\).
\end{theorem}
\begin{proofbox}
\emph{Invariant.} Induction over DFS events.  A child region \(C\) merges into
\(u\)'s side exactly when \(C\) is maximal-BLACK (child returned) and \(u\) is
still GREY, so sets are maximal BLACK regions glued under their waiting GREY
parent; stamping \(\mathit{anc}[\mathit{find}(u)]=u\) after each merge keeps the
stamp that parent. \emph{Correctness.} Suppose \(u\) turns BLACK and partner
\(v\) is already BLACK (wlog \(v\) died first, or \(v\) is in an earlier branch).
At \(u\)'s death moment, \(\mathit{find}(v)\)'s region is the maximal subtree
rooted at some GREY node's completed child-chain, and its stamp \(w\) is the
deepest GREY ancestor of \(v\).  GREY nodes form exactly the current root-to-\(u\)
stack, so \(w\) is on that stack.  \(w\) ancestors \(v\) and \(w\) ancestors
\(u\) (it is the DFS-stack node currently processing the region that contains
\(u\)); any deeper stack node \(z\) child-of-\(w\) toward \(u\) cannot contain
\(v\) (if it did, \(v\)'s BLACK region would have merged up through \(z\) into
\(w\) only after \(z\)'s return -- and \(z\) is GREY now, with \(u\) inside it,
while \(v\) died inside an \emph{already returned} child of \(w\)).  Hence
\(w=\lca(u,v)\). \emph{Cost.} \(n-1\) unions, \(O(n+q)\) finds, DSU with path
compression (union here always points a child set at its parent set, so even
without rank the amortised bound is \(\bigO{\alpha}\) for the mixed batch) --
total \(\bigO{(n+q)\alpha(n)}\). \(\Box\)
\end{proofbox}

\begin{lstlisting}[caption={Tarjan offline LCA in full: two-endpoint query attachment, union-on-child-return, answer on second endpoint's death. ASCII only}, label={lst:c21-tarjan}]
#include <bits/stdc++.h>
using namespace std;
const int MAXN=200005;
vector<int> adj[MAXN];
vector<pair<int,int>> qlist[MAXN];      // (partner, query id), on BOTH ends
int fa[MAXN], anc[MAXN], col[MAXN];     // col: 0 WHITE, 1 GREY, 2 BLACK
int ans[MAXN];
int find(int x){ return fa[x]==x ? x : fa[x]=find(fa[x]); }
void dfs(int u, int p){
  col[u] = 1;  fa[u] = u;  anc[u] = u;
  for (int v : adj[u]) if (v != p){
    dfs(v, u);
    fa[find(v)] = find(u);              // finished region joins the parent
    anc[find(u)] = u;                   // ... whose stamp is u
  }
  for (auto &Q : qlist[u])
    if (col[Q.first] == 2)              // partner already dead -> answer now
      ans[Q.second] = anc[find(Q.first)];
  col[u] = 2;
}
int main(){
  // read n, q, edges; attach: qlist[u].push_back({v,i}); qlist[v].push_back({u,i});
  // dfs(1, 0);  print ans[0..q-1]
}
\end{lstlisting}

% ---------------------------------------------------------------
\subsection{Worked example: five queries on \(R\), every event traced}
% ---------------------------------------------------------------

Run \(\coderef{c21-tarjan}\) on the running tree \(R\) (\(\figref{c21-tree}\)),
children visited left-to-right, DFS preorder \(P=[1,2,4,8,5,3,6,7]\), with the
batch:
\[
  Q_1=(8,6),\quad Q_2=(4,5),\quad Q_3=(2,8),\quad Q_4=(6,7),\quad Q_5=(5,7).
\]
Expected answers (Ch.~16 by inspection of the picture):
\(\lca(8,6){=}1\), \(\lca(4,5){=}2\), \(\lca(2,8){=}2\), \(\lca(6,7){=}3\),
\(\lca(5,7){=}1\).

\begin{tracebox}[The full event log -- unions, stamps, and the five answer moments]
\textbf{Enter} 1,2,4,8 (GREY stack \([1,2,4,8]\)). \
8: no children; its lists \(Q_1,Q_3\): partners 6 WHITE, 2 GREY -- \emph{no
answer}; \textbf{BLACK}(8). \
Return 8\(\to\)4: \(\mathit{union}(4,8)\), set \(\{4,8\}\), stamp \(4\). \
4: lists \(Q_2,Q_3\): partner 5 WHITE, 2 GREY -- no answer;
\textbf{BLACK}(4). \
Return 4\(\to\)2: \(\mathit{union}(2,4)\), set \(\{2,4,8\}\), stamp \(2\). \
Enter 5; 5 dies: \(Q_2\): partner 4 \textbf{BLACK} \(\Rightarrow\)
\(\mathit{ans}_2=\mathit{anc}[\mathit{find}(4)]=\boxed{2}\) ✓;
\(Q_5\): partner 7 WHITE -- wait; \textbf{BLACK}(5). \
Return 5\(\to\)2: \(\mathit{union}(2,5)\), set \(\{2,4,5,8\}\), stamp \(2\). \
2 dies: \(Q_3\): partner 8 BLACK \(\Rightarrow\)
\(\mathit{ans}_3=\mathit{anc}[\mathit{find}(8)]=\boxed{2}\) ✓ (note 8's set is
\(\{2,4,5,8\}\) stamped \(2\) -- the ancestor case answered one union after
8's own death). \
Return 2\(\to\)1: \(\mathit{union}(1,2)\), set \(\{1,2,4,5,8\}\), stamp \(1\). \
Enter 3,6; 6 dies: \(Q_1\): partner 8 BLACK \(\Rightarrow\)
\(\mathit{ans}_1=\mathit{anc}[\mathit{find}(8)]=\boxed{1}\) ✓ -- the two
endpoints' regions met under 1; \(Q_4\): 7 WHITE -- wait. \
Return 6\(\to\)3: \(\mathit{union}(3,6)\), set \(\{3,6\}\), stamp \(3\). \
Enter 7; 7 dies: \(Q_4\): partner 6 BLACK \(\Rightarrow\)
\(\mathit{ans}_4=\mathit{anc}[\mathit{find}(6)]=\boxed{3}\) ✓ (siblings meet
under 3); \(Q_5\): partner 5 BLACK \(\Rightarrow\)
\(\mathit{ans}_5=\mathit{anc}[\mathit{find}(5)]=\boxed{1}\) ✓ (5's set merged
into 1 at 2's death). \
Return 7\(\to\)3: \(\mathit{union}(3,7)\); 3 dies (empty lists); return,
\(\mathit{union}(1,3)\); 1 dies.  Answers: \(1,2,2,3,1\) -- all five match.
\(\Box\)
\end{tracebox}

\begin{insightbox}[Read the invariant off the trace in one line]
At every answer moment, \(\mathit{find}(\text{partner})\)'s set was "everything
already finished in one branch", and the stamp was "the still-open node
directly above it" -- which, because GREY nodes are exactly the current
root-to-\(u\) stack, is precisely where the two endpoints' paths merge.  Tarjan
is not \emph{computing} ancestors; it is watching the DFS stack and naming the
merge point at the instant both sides stop moving.
\end{insightbox}

\begin{comparisonbox}[Tarjan vs.\ lifting (Ch.~16) vs.\ Euler+RMQ (Ch.~5)]
\begin{itemize}[nosep,leftmargin=1.4em]
  \item \textbf{Modality.} Tarjan: pure offline -- a query asked after an
        update, or answered out of batch order, is outside its contract.
        Lifting (\(\chref{c11}\)): online, \(O(\log n)\)/query, table
        \(\bigO{n\log n}\).  Euler tour + RMQ: offline-or-online after one
        \(\bigO{n}\) build plus \(\bigO{n\log n}\) (or linear Cartesian) RMQ.
  \item \textbf{Cost at \(q \approx n\) large:} Tarjan \(\bigO{(n+q)\alpha}\)
        \(<\) both \(\log\) methods in theory, and \emph{much} less memory --
        no table.  This is why contest mega-batches (SPOJ LCA, problem data with
        \(n,q\ge10^5\) and tight memory) hand themselves to it.
  \item \textbf{Fragility.} Tarjan needs all queries upfront \emph{and} a
        recursion-depth-safe DFS at \(n=2\cdot10^5\); lifting needs neither but
        pays the table.  Path Mo (\(\secref{c21-path}\)) needs \(\lca\) only to
        \emph{sort windows}, so batch its LCAs through Tarjan and the whole
        pipeline stays \(\bigO{(n+q)\sqrt{n+q}}\) with one \(\bigO{n}\) prep.
\end{itemize}
\end{comparisonbox}

\begin{complexitybox}[Tarjan]
One DFS: \(n-1\) unions; each query touched twice on attach, one find at the
answering endpoint (the WHITE/GREY partner check on the other endpoint costs
\(O(1)\)).  With path compression the DSU batch is \(\bigO{(n+q)\alpha(n)}\)
\(\approx\) linear.  \textbf{Space} \(\bigO{n+q}\) -- \emph{no lifting table},
which for \(n=10^6\) is the difference between fitting and not fitting.  The
recursion is the real risk: convert to an explicit stack when \(n\) is large.
\end{complexitybox}

\begin{warningbox}[Three ways Tarjan goes wrong]
\begin{enumerate}[nosep,leftmargin=1.6em]
  \item \textbf{Attaching the query to only one endpoint.} The answer fires at
        whichever endpoint turns BLACK \emph{second}; if the query sits only on
        the first one, it checks a WHITE/GREY partner, skips, and the query is
        silently never answered (stale \(0\) in \(\mathit{ans}\)).
  \item \textbf{Stamping after \(\mathit{find}(v)\) instead of \(\mathit{find}(u)\)}
        (or forgetting the stamp entirely): the set survives but its
        \(\mathit{anc}\) entry is wrong, and answers drift up or down by one
        level -- test on the \(Q_3=(2,8)\) ancestor shape, which dies first to
        this bug.
  \item \textbf{Checking \(\mathit{col}[\text{partner}] != \mathit{WHITE}\) instead
        of \(=\)BLACK:} a GREY partner means "query is being asked inside the
        partner's unfinished subtree"; answering from its stamp at that moment
        can look right (\(\mathit{anc}[\mathit{find}(\text{partner})]\) is the
        partner or its child) only because of an accident of that variant's
        invariants -- the clean contract is: \emph{answer exactly when the
        partner is BLACK}.  (The one-flag implementations you see on the web
        fold "GREY ancestor" into BLACK's formula -- derive it yourself before
        trusting it.)
\end{enumerate}
\end{warningbox}

\begin{hardeningbox}
\begin{itemize}[nosep,leftmargin=1.4em]
  \item \textbf{Same skeleton, other offline reads.} The union-on-return +
        death-time-scan pattern computes more than LCA: batch "nearest marked
        ancestor", offline subtree counts at death time -- the scan list is a
        general "what is visible from the stack right now" hook.
  \item \textbf{Iterative form.} Push \((u,\text{iterator state})\) frames; the
        union and the scan both happen on \emph{pop} -- one loop, no stack
        overflow at \(n=2\cdot10^5\).
  \item \textbf{Forest input?} Loop \(\mathit{dfs}\) over WHITE roots; queries
        across components must be answered separately (no LCA) -- check whether
        the statement guarantees connectivity before trusting \(\mathit{ans}\).
\end{itemize}
\end{hardeningbox}

\begin{recallbox}[27.4 in one breath]
\emph{One DFS, three colours; union a returned child into the parent and stamp
the parent; a query fires when its second endpoint turns BLACK, reading
\(\mathit{anc}[\mathit{find}(\text{partner})]\) -- the merge point of the two
endpoints' paths (\(\defref{c21-tarjan}\), \(\thmref{c21-tarjan}\),
\(\coderef{c21-tarjan}\)).  The five-query log on \(R\) returns
\(1,2,2,3,1\).  Batch modality \(\Rightarrow\) Tarjan; online \(\Rightarrow\)
Ch.~16 lifting.}
\selftest{
  \item What is \(\mathit{anc}[\mathit{find}(x)]\) mid-DFS? \emph{(The deepest GREY node whose subtree contains the finished region \(x\)'s set belongs to -- the invariant of \thmref{c21-tarjan}.)}
  \item Why attach queries to both endpoints? \emph{(The firing moment is the second endpoint's death; either endpoint can be second.)}
  \item State the cost difference against Ch.~16 at \(n=q=10^6\), memory-wise. \emph{(Tarjan O((n+q)alpha), space O(n+q), no table; lifting O((n+q)log n) with an O(n log n) table -- often the memory ceiling.)}
}
\end{recallbox}
